% ==================================================
%  CAPÍTULO 1: INTRODUCCIÓN AL DEEP LEARNING
% ==================================================

\section{Introducción al Deep Learning}

\subsection{Contexto tecnológico: la Sociedad 5.0}

En 2015, Japón acuñó el concepto de \textbf{Sociedad 5.0}, cuya idea central es que \emph{las máquinas serán las principales creadoras del conocimiento humano}, apalancadas en la inteligencia artificial. Mientras la Sociedad 4.0 generó la información digital, la 5.0 la aprovecha con IA para integrar la tecnología en la vida cotidiana.

\subsection{Inteligencia Artificial: definición y evolución}

Se adopta la definición de Nilsson (1998):

\begin{definicion}
\textbf{Inteligencia Artificial (IA)}: capacidad de imitar comportamientos inteligentes.
\end{definicion}

La definición de IA ha evolucionado a lo largo del tiempo.

\begin{table}[H]
\centering
\caption{Evolución de la definición de Inteligencia Artificial.}
\label{tab:evolucion-ia}
\small
\begin{tabularx}{\textwidth}{@{}lXX@{}}
\toprule
\textbf{Definición histórica} & \textbf{Idea central} & \textbf{Problema / observación} \\
\midrule
Sistemas que piensan como humanos & Imitar el pensamiento humano & No existe una única forma de pensamiento humano \\
Sistemas que piensan racionalmente & Pensar de manera correcta & El silogismo carece de contexto \\
Sistemas que actúan como humanos & Test de Turing & No tuvo éxito en su época por capacidad de procesamiento \\
Sistemas que actúan racionalmente & Actuar según lo programado & Definición vigente desde 1998 \\
\bottomrule
\end{tabularx}
\end{table}

\begin{definicion}
\textbf{IA (definición vigente)}: todo sistema capaz de \emph{actuar racionalmente}.

\textbf{Racional}: correcto para lo que fue programado; no implica conciencia.
\end{definicion}

\subsubsection{Clasificaciones de la Inteligencia Artificial}

\begin{table}[H]
\centering
\caption{Clasificación de la IA según capacidad y complejidad.}
\label{tab:ia-capacidad}
\small
\begin{tabularx}{\textwidth}{@{}lXX@{}}
\toprule
\textbf{Tipo} & \textbf{Descripción} & \textbf{Ejemplos} \\
\midrule
IA reactiva & No aprende de experiencias pasadas; responde a estímulos actuales & Alexa, Siri, aspiradora robot, Deep Blue \\
IA con memoria limitada & Aprende de datos históricos, pero con memoria limitada y específica & Vehículos autónomos, Waze, Google Maps \\
\bottomrule
\end{tabularx}
\end{table}

\begin{table}[H]
\centering
\caption{Clasificación de la IA según propósito y aplicaciones.}
\label{tab:ia-proposito}
\small
\begin{tabularx}{\textwidth}{@{}lXX@{}}
\toprule
\textbf{Tipo} & \textbf{Descripción} & \textbf{Ejemplos} \\
\midrule
IA débil / narrow & Específica para tareas claras; no aprende tareas nuevas & Detector de spam de Gmail, asistentes virtuales, sistemas de recomendación, reconocimiento de voz \\
IA fuerte & Capaz de aprender cosas nuevas & Autos autónomos, sistemas de Machine Learning \\
Super IA & Consciente de su existencia; capaz de realizar cualquier tarea humana & No existe todavía \\
\bottomrule
\end{tabularx}
\end{table}

\subsection{Machine Learning: el núcleo de la IA}

\begin{definicion}
\textbf{Machine Learning} (aprendizaje automático): campo de la IA que se encarga de reconocer
patrones desde los datos y de generalizar conocimiento a partir de la experiencia y de muchos
datos.

Definición de \textbf{Arthur Samuel}: ``el diseño y desarrollo de algoritmos que permiten a los
computadores \emph{aprender sin ser explícitamente programados}''.
\end{definicion}

No existe un comando ``\texttt{.learn}'', el aprendizaje se logra mediante un \textbf{algoritmo}.

\begin{definicion}
\textbf{Algoritmo}: secuencia lógica, paso a paso, con inicio, decisiones y fin u objetivo.
\end{definicion}

\begin{table}[H]
\centering
\caption{Tipos de aprendizaje automático.}
\label{tab:tipos-ml}
\small
\begin{tabularx}{\textwidth}{@{}lXX@{}}
\toprule
\textbf{Tipo} & \textbf{Descripción} & \textbf{Ejemplos} \\
\midrule
Supervisado & Aprende con ejemplos etiquetados & Detección de spam, reconocimiento de voz, analíticas de anuncios \\
No supervisado & Encuentra patrones sin etiquetas previas & Sistemas de recomendación, registros de conductas, detección de anomalías \\
Semi-supervisado & Mezcla de supervisado y no supervisado & Análisis de voz, detección de fraude, investigaciones médicas \\
Por refuerzo & Aprende mediante prueba, error y refuerzo & Robótica, videojuegos, gestión de recursos, IoT \\
\bottomrule
\end{tabularx}
\end{table}

\subsection{Deep Learning: aprendizaje por capas}

\begin{definicion}
\textbf{Deep Learning} (aprendizaje profundo): sistema capaz de aprender mediante algoritmos y
\emph{sin intervención manual} de un usuario.
\end{definicion}

\begin{definicion}
\textbf{Representación}: forma en que los datos son transformados y procesados a lo largo de
las capas para extraer características relevantes de manera automática.

\textbf{Profundidad del modelo}: cantidad de capas sucesivas de representación.
\end{definicion}

\begin{advertencia}
``Profundo'' \emph{no} significa comprensión más profunda; significa \emph{muchas capas
sucesivas de representación}.
\end{advertencia}

\textbf{¿Por qué despegó después de 2012?}

Las ideas clave de la visión artificial (redes neuronales convolucionales y retropropagación)
ya se entendían en 1990 y el algoritmo \textbf{LSTM} (memoria a corto y largo plazo) se
desarrolló en 1997 y apenas ha cambiado. Pero hasta 2010, un computador de un terabyte era muy difícil
de conseguir.

\textbf{Tres nuevas figuras han impulsado el avance del aprendizaje automático:}
\begin{enumerate}[topsep=2pt,itemsep=1pt]
  \item Hardware más potente (GPU, TPU).
  \item Mayor disponibilidad de datos (big data).
  \item Avances en los algoritmos (mejores funciones de activación, optimizadores).
\end{enumerate}

\subsubsection{Machine Learning vs Deep Learning}

\begin{table}[H]
\centering
\caption{Comparación entre Machine Learning y Deep Learning.}
\label{tab:ml-vs-dl}
\small
\begin{tabularx}{\textwidth}{@{}lXX@{}}
\toprule
\textbf{Criterio} & \textbf{Machine Learning} & \textbf{Deep Learning} \\
\midrule
Extracción de variables & Manual, mediante ingeniería de características & Automática, mediante capas de representación \\
Datos & Principalmente estructurados & Imágenes, texto, voz y otros no estructurados \\
Intervención manual & Requiere decirle si aprendió bien o mal & Aprende sin intervención manual \\
Escalabilidad & Computacionalmente eficiente, pero no se adapta bien a grandes volúmenes & Suele ser más eficaz en problemas complejos y grandes volúmenes \\
Costo computacional & Menor & Mayor \\
Interpretabilidad & Mayor & Menor \\
\bottomrule
\end{tabularx}
\end{table}

\begin{advertencia}
El Deep Learning es más caro y menos interpretable, pero más eficaz en problemas complejos.
\end{advertencia}